\section{All factor pairs above one integer}
\label{sec:full-fibre}
%======================================================================

A fixed integer has many factor pairs.  This section describes the complete geometric fibre first, so that the later canonical M\"obius point is not confused with the whole factorization geometry.

The factorization geometry and the smooth--transport geometry use the same coordinate map but different point selections.  This distinction is essential.

Fix $N\ge1$ and consider every positive factor pair
\begin{equation}
 uv=N,
 \qquad
 1\le u\le v.
 \label{eq:factor-pair}
\end{equation}
Define the Fermat coordinates
\begin{equation}
 a=\frac{u+v}{2},
 \qquad
 b=\frac{v-u}{2}\ge0.
 \label{eq:full-Fermat}
\end{equation}
Then
\begin{equation}
 u=a-b,
 \qquad
 v=a+b,
 \qquad
 N=a^2-b^2.
 \label{eq:full-recovery}
\end{equation}
The coordinates lie in $\frac12\Z$; they are integers when $u$ and $v$ have the same parity.

Squaring the complex factor gives
\begin{equation}
 (a+bi)^2=N+iY,
 \qquad
 \boxed{Y=2ab=\frac{v^2-u^2}{2}.}
 \label{eq:full-square}
\end{equation}
Thus all factorizations of a fixed $N$ lie on the vertical fibre
\begin{equation}
 \Repart w=N.
 \label{eq:vertical-fibre}
\end{equation}

\subsection{Square-root tests}

Let
\begin{equation}
 w=N+iY,
 \qquad
 \rho=|w|=\sqrt{N^2+Y^2}.
 \label{eq:rho-full}
\end{equation}
Then
\begin{align}
 a^2&=\frac{\rho+N}{2},
 \label{eq:a-square}\\
 b^2&=\frac{\rho-N}{2},
 \label{eq:b-square}\\
 u^2&=\rho-Y,
 \label{eq:u-square}\\
 v^2&=\rho+Y.
 \label{eq:v-square}
\end{align}
Accordingly, a candidate vertical ordinate $Y$ corresponds to an exact factorization if and only if the relevant square and parity conditions hold.  The identity
\begin{equation}
 \sqrt{N^2+Y^2}=a^2+b^2
 \label{eq:hypotenuse}
\end{equation}
is the right-triangle form of the same condition.

\subsection{Admissible intervals}

For a factor pair $u(N/u)=N$ with $1\le u\le\sqrt N$,
\begin{equation}
 a(u)=\frac{u+N/u}{2},
 \qquad
 b(u)=\frac{N/u-u}{2}.
 \label{eq:ab-u}
\end{equation}
Both functions decrease as $u$ increases toward $\sqrt N$.  Therefore the full positive factor fibre satisfies
\begin{align}
 \sqrt N&\le a\le\frac{N+1}{2},
 \label{eq:a-full-range}\\
 0&\le b\le\frac{N-1}{2},
 \label{eq:b-full-range}\\
 0&\le Y\le\frac{N^2-1}{2}.
 \label{eq:Y-full-range}
\end{align}
For integer enumeration one uses $a\ge\lceil\sqrt N\rceil$ with the parity required by $N$.

The upper endpoint corresponds to the trivial factor pair $1\cdot N$.  If the trivial factor is excluded and the smallest permitted odd factor is $3$, then
\begin{align}
 a&\le\frac{N+9}{6},
 \label{eq:a-three}\\
 b&\le\frac{N-9}{6},
 \label{eq:b-three}\\
 Y&\le\frac{N^2-81}{18}.
 \label{eq:Y-three}
\end{align}
These are the nontrivial odd-factor or semiprime-search envelopes.  They are not the universal bounds for the complete factor fibre.

\begin{example}[The fibre $N=309$]
The trivial factor pair $1\cdot309$ gives
\[
 (a,b)=(155,154),
 \qquad
 w=309+47740i.
\]
The nontrivial factorization $3\cdot103$ gives
\[
 (a,b)=(53,50),
 \qquad
 w=309+5300i.
\]
The latter is the highest point allowed by the nontrivial odd-factor condition $u\ge3$, but not the highest point on the full factor fibre.
\end{example}

\begin{figure}[htbp]
\centering
\begin{tikzpicture}[>=Latex,node distance=13mm]
  \node[draw,rounded corners,fill=blue!5,align=center,minimum width=31mm,minimum height=12mm] (factor) {$uv=N$\\factor coordinates};
  \node[draw,rounded corners,fill=green!6,align=center,minimum width=35mm,minimum height=12mm,right=of factor] (fermat) {$a=(u+v)/2$\\$b=(v-u)/2$};
  \node[draw,rounded corners,fill=red!5,align=center,minimum width=38mm,minimum height=12mm,right=of fermat] (square) {$w=N+iY$\\$Y=2ab$};
  \node[draw,rounded corners,fill=gray!8,align=center,minimum width=38mm,minimum height=12mm,below=of square] (null) {$U=|w|-Y=u^2$\\$V=|w|+Y=v^2$};
  \draw[->,thick] (factor) -- (fermat);
  \draw[->,thick] (fermat) -- node[above,font=\small] {$z\mapsto z^2$} (square);
  \draw[->,thick] (square) -- (null);
  \draw[->,thick,bend left=20] (null.west) to node[below,font=\small] {exact recovery} (factor.south);
\end{tikzpicture}
\caption{Equivalent coordinate systems on a complete factor fibre.  The squared coordinates retain the factor pair exactly.}
\label{fig:coordinate-map}
\end{figure}

%======================================================================
